\documentclass[10pt,a4paper]{amsart}
\usepackage{amsmath,amssymb,amsthm,graphicx}
\usepackage{verbatim}
\newenvironment{reference}{\comment}{\endcomment}
%\newenvironment{reference}{}{}
\newenvironment{slogan}{\comment}{\endcomment}
\newenvironment{history}{\comment}{\endcomment}

% For commutative diagrams we use Xy-pic
\usepackage[all]{xy}

% We use 2cell for 2-commutative diagrams.
\xyoption{2cell}
% Theorem environments.
%
\theoremstyle{plain}
\newtheorem{theorem}[subsection]{Theorem}
\newtheorem{proposition}[subsection]{Proposition}
\newtheorem{lemma}[subsection]{Lemma}

\theoremstyle{definition}
\newtheorem{definition}[subsection]{Definition}
\newtheorem{example}[subsection]{Example}
\newtheorem{exercise}[subsection]{Exercise}
\newtheorem{situation}[subsection]{Situation}

\theoremstyle{remark}
\newtheorem{remark}[subsection]{Remark}
\newtheorem{remarks}[subsection]{Remarks}

\numberwithin{equation}{subsection}

% Macros
%
\def\lim{\mathop{\mathrm{lim}}\nolimits}
\def\colim{\mathop{\mathrm{colim}}\nolimits}
\def\Spec{\mathop{\mathrm{Spec}}}
\def\Hom{\mathop{\mathrm{Hom}}\nolimits}
\def\Ext{\mathop{\mathrm{Ext}}\nolimits}
\def\SheafHom{\mathop{\mathcal{H}\!\mathit{om}}\nolimits}
\def\SheafExt{\mathop{\mathcal{E}\!\mathit{xt}}\nolimits}
\def\Sch{\mathit{Sch}}
\def\Mor{\mathop{\mathrm{Mor}}\nolimits}
\def\Ob{\mathop{\mathrm{Ob}}\nolimits}
\def\Sh{\mathop{\mathit{Sh}}\nolimits}
\def\NL{\mathop{N\!L}\nolimits}
\def\CH{\mathop{\mathrm{CH}}\nolimits}
\def\proetale{{pro\text{-}\acute{e}tale}}
\def\etale{{\acute{e}tale}}
\def\QCoh{\mathit{QCoh}}
\def\Ker{\mathop{\mathrm{Ker}}}
\def\Im{\mathop{\mathrm{Im}}}
\def\Coker{\mathop{\mathrm{Coker}}}
\def\Coim{\mathop{\mathrm{Coim}}}

% Boxtimes
%
\DeclareMathSymbol{\boxtimes}{\mathbin}{AMSa}{"02}

%
% Macros for moduli stacks/spaces
%
\def\QCohstack{\mathcal{QC}\!\mathit{oh}}
\def\Cohstack{\mathcal{C}\!\mathit{oh}}
\def\Spacesstack{\mathcal{S}\!\mathit{paces}}
\def\Quotfunctor{\mathrm{Quot}}
\def\Hilbfunctor{\mathrm{Hilb}}
\def\Curvesstack{\mathcal{C}\!\mathit{urves}}
\def\Polarizedstack{\mathcal{P}\!\mathit{olarized}}
\def\Complexesstack{\mathcal{C}\!\mathit{omplexes}}
% \Pic is the operator that assigns to X its picard group, usage \Pic(X)
% \Picardstack_{X/B} denotes the Picard stack of X over B
% \Picardfunctor_{X/B} denotes the Picard functor of X over B
\def\Pic{\mathop{\mathrm{Pic}}\nolimits}
\def\Picardstack{\mathcal{P}\!\mathit{ic}}
\def\Picardfunctor{\mathrm{Pic}}
\def\Deformationcategory{\mathcal{D}\!\mathit{ef}}
\usepackage{cleveref}
\crefname{originallemma}{original lemma}{original lemmas}
\Crefname{originallemma}{Original lemma}{Original lemmas}
\makeatletter
\newlabel{descent-lemma-descending-property-finite-presentation}{{02L0}{source}{}{sourcepointer-1}{}}
\newlabel{descent-lemma-descending-property-finite-presentation@cref}{{[originallemma][0][]02L0}{[1][0][]source}{}{}{}}
\makeatother
\crefname{originallemma}{original lemma}{original lemmas}
\Crefname{originallemma}{Original lemma}{Original lemmas}
\makeatletter
\newlabel{descent-lemma-descending-property-monomorphism}{{02YK}{source}{}{sourcepointer-2}{}}
\newlabel{descent-lemma-descending-property-monomorphism@cref}{{[originallemma][0][]02YK}{[1][0][]source}{}{}{}}
\makeatother
\crefname{originallemma}{original lemma}{original lemmas}
\Crefname{originallemma}{Original lemma}{Original lemmas}
\makeatletter
\newlabel{descent-lemma-descent-data-sheaves}{{02W5}{source}{}{sourcepointer-3}{}}
\newlabel{descent-lemma-descent-data-sheaves@cref}{{[originallemma][0][]02W5}{[1][0][]source}{}{}{}}
\makeatother
\crefname{originallemma}{original lemma}{original lemmas}
\Crefname{originallemma}{Original lemma}{Original lemmas}
\makeatletter
\newlabel{flat-lemma-flat-dimension-n-representable}{{05UF}{source}{}{sourcepointer-4}{}}
\newlabel{flat-lemma-flat-dimension-n-representable@cref}{{[originallemma][0][]05UF}{[1][0][]source}{}{}{}}
\makeatother
\crefname{originalreference}{original reference}{original references}
\Crefname{originalreference}{Original reference}{Original references}
\makeatletter
\newlabel{lemma-existence-complete}{{05HU}{source}{}{sourcepointer-5}{}}
\newlabel{lemma-existence-complete@cref}{{[originalreference][0][]05HU}{[1][0][]source}{}{}{}}
\makeatother
\crefname{originalreference}{original reference}{original references}
\Crefname{originalreference}{Original reference}{Original references}
\makeatletter
\newlabel{lemma-flat-dimension-n}{{05MV}{source}{}{sourcepointer-6}{}}
\newlabel{lemma-flat-dimension-n@cref}{{[originalreference][0][]05MV}{[1][0][]source}{}{}{}}
\makeatother
\crefname{originalreference}{original reference}{original references}
\Crefname{originalreference}{Original reference}{Original references}
\makeatletter
\newlabel{lemma-localize-flat-dimension-n}{{05UD}{source}{}{sourcepointer-7}{}}
\newlabel{lemma-localize-flat-dimension-n@cref}{{[originalreference][0][]05UD}{[1][0][]source}{}{}{}}
\makeatother
\crefname{originalreference}{original reference}{original references}
\Crefname{originalreference}{Original reference}{Original references}
\makeatletter
\newlabel{lemma-pre-flat-dimension-n}{{0CWF}{source}{}{sourcepointer-8}{}}
\newlabel{lemma-pre-flat-dimension-n@cref}{{[originalreference][0][]0CWF}{[1][0][]source}{}{}{}}
\makeatother
\crefname{originalreference}{original reference}{original references}
\Crefname{originalreference}{Original reference}{Original references}
\makeatletter
\newlabel{lemma-quasi-finite-base-change}{{00PP}{source}{}{sourcepointer-9}{}}
\newlabel{lemma-quasi-finite-base-change@cref}{{[originalreference][0][]00PP}{[1][0][]source}{}{}{}}
\makeatother
\crefname{originallemma}{original lemma}{original lemmas}
\Crefname{originallemma}{Original lemma}{Original lemmas}
\makeatletter
\newlabel{more-morphisms-lemma-separated-locally-quasi-finite-morphisms-fppf-descend}{{02W8}{source}{}{sourcepointer-10}{}}
\newlabel{more-morphisms-lemma-separated-locally-quasi-finite-morphisms-fppf-descend@cref}{{[originallemma][0][]02W8}{[1][0][]source}{}{}{}}
\makeatother
\crefname{originallemma}{original lemma}{original lemmas}
\Crefname{originallemma}{Original lemma}{Original lemmas}
\makeatletter
\newlabel{morphisms-lemma-constructible-containing-open}{{05LW}{source}{}{sourcepointer-11}{}}
\newlabel{morphisms-lemma-constructible-containing-open@cref}{{[originallemma][0][]05LW}{[1][0][]source}{}{}{}}
\makeatother
\crefname{originallemma}{original lemma}{original lemmas}
\Crefname{originallemma}{Original lemma}{Original lemmas}
\makeatletter
\newlabel{morphisms-lemma-monomorphism-loc-finite-type-loc-quasi-finite}{{0CT8}{source}{}{sourcepointer-12}{}}
\newlabel{morphisms-lemma-monomorphism-loc-finite-type-loc-quasi-finite@cref}{{[originallemma][0][]0CT8}{[1][0][]source}{}{}{}}
\makeatother
\crefname{originalreference}{original reference}{original references}
\Crefname{originalreference}{Original reference}{Original references}
\makeatletter
\newlabel{situation-flat-dimension-n}{{05MT}{source}{}{sourcepointer-13}{}}
\newlabel{situation-flat-dimension-n@cref}{{[originalreference][0][]05MT}{[1][0][]source}{}{}{}}
\makeatother
\crefname{originallemma}{original lemma}{original lemmas}
\Crefname{originallemma}{Original lemma}{Original lemmas}
\makeatletter
\newlabel{spaces-divisors-lemma-relative-assassin-constructible}{{0CVV}{source}{}{sourcepointer-14}{}}
\newlabel{spaces-divisors-lemma-relative-assassin-constructible@cref}{{[originallemma][0][]0CVV}{[1][0][]source}{}{}{}}
\makeatother
\title{The locus where a finite presentation module is flat in dimensions at least n is representable by a monomorphism of finite presentation}
\author{The Stacks Project Authors}
\date{stacks-project:spaces-flat.tex:local-snapshot-20261012; GNU FDL 1.2 or later}
\begin{document}
\maketitle
\noindent\textit{Editorial note.} All delivered mathematical text is the author's own, copied from the pinned Stacks Project chapter file spaces-flat.tex; no statement or proof step was rewritten, reordered or removed.\par\medskip
\noindent\textit{Editorial note.} The counted proof is the author's complete argument: a reduction of representability to a single etale cover, a descending induction on n through quasi-finite base change, and a complete devissage plus constructibility argument reducing to the flat case\par\medskip
\noindent\textit{Editorial note.} The label theorem-flat-dimension-n-representable names the counted statement, and the proof environment immediately following it in the source is delivered verbatim as its proof.\par\medskip
\noindent\textit{Editorial note.} The counted statement is carried by the author's own theorem environment in the chapter source.\par\medskip
\noindent\textit{Editorial note.} The counted statement is the author's theorem environment labelled theorem-flat-dimension-n-representable (2195-2203) and the counted proof is the proof environment beginning at 2205, so the pairing is by adjacency.\par\medskip
\noindent\textit{Editorial note.} The proof is complete as the author wrote it: every reduction is justified where it is made and the only external input is the flat case quoted at the end, which is recorded among the allowed prerequisites.\par\medskip
\noindent\textit{Editorial note.} All delivered mathematical text is the author's own, copied from the pinned Stacks Project chapter file spaces-flat.tex; no statement, proof step or symbol was rewritten, reordered or removed.\par\medskip
\noindent\textit{Allowed standard prerequisites (editorial).}\begin{itemize}
\item The functor classifying flatness in dimensions at least n, its sheaf property for the fppf topology, and the fact that it is a subfunctor of the corresponding functor with a larger bound
\item Purity of a module relative to the base, its preservation under quasi-finite base change and under etale localization, and the constructibility of the relative assassin
\item Etale descent for monomorphisms and for morphisms of finite presentation, the existence of a complete devissage in an etale neighbourhood, and the flat case of the statement
\end{itemize}\medskip
\ifdefined\hypertarget\hypertarget{sourcepointer-1}{}\fi
\noindent\textit{Source reference.} Original lemma 02L0: sources/stacks-chow/tags.txt, source lines 3391--3391. Prerequisite pointer invoked at proof line(s) 2240; it names a standard fact of a sibling chapter that the counted argument uses, and every step of the counted proof built on it is retained verbatim in the delivered body.\par\medskip
\ifdefined\hypertarget\hypertarget{sourcepointer-2}{}\fi
\noindent\textit{Source reference.} Original lemma 02YK: sources/stacks-chow/tags.txt, source lines 3879--3879. Prerequisite pointer invoked at proof line(s) 2239; it names a standard fact of a sibling chapter that the counted argument uses, and every step of the counted proof built on it is retained verbatim in the delivered body.\par\medskip
\ifdefined\hypertarget\hypertarget{sourcepointer-3}{}\fi
\noindent\textit{Source reference.} Original lemma 02W5: sources/stacks-chow/tags.txt, source lines 3792--3792. Prerequisite pointer invoked at proof line(s) 2235; it names a standard fact of a sibling chapter that the counted argument uses, and every step of the counted proof built on it is retained verbatim in the delivered body.\par\medskip
\ifdefined\hypertarget\hypertarget{sourcepointer-4}{}\fi
\noindent\textit{Source reference.} Original lemma 05UF: sources/stacks-chow/tags.txt, source lines 7514--7514. Prerequisite pointer invoked at proof line(s) 2307; it names a standard fact of a sibling chapter that the counted argument uses, and every step of the counted proof built on it is retained verbatim in the delivered body.\par\medskip
\ifdefined\hypertarget\hypertarget{sourcepointer-5}{}\fi
\noindent\textit{Source reference.} Original reference 05HU: sources/stacks-spaces-flat/pinned-original-source.tex, source lines 640--640. Prerequisite pointer invoked at proof line(s) 2272; it names a standard fact of the same chapter that the counted argument uses, and every step of the counted proof built on it is retained verbatim in the delivered body.\par\medskip
\ifdefined\hypertarget\hypertarget{sourcepointer-6}{}\fi
\noindent\textit{Source reference.} Original reference 05MV: sources/stacks-spaces-flat/pinned-original-source.tex, source lines 2122--2122. Prerequisite pointer invoked at proof line(s) 2207; it names a standard fact of the same chapter that the counted argument uses, and every step of the counted proof built on it is retained verbatim in the delivered body.\par\medskip
\ifdefined\hypertarget\hypertarget{sourcepointer-7}{}\fi
\noindent\textit{Source reference.} Original reference 05UD: sources/stacks-spaces-flat/pinned-original-source.tex, source lines 2169--2169. Prerequisite pointer invoked at proof line(s) 2296; it names a standard fact of the same chapter that the counted argument uses, and every step of the counted proof built on it is retained verbatim in the delivered body.\par\medskip
\ifdefined\hypertarget\hypertarget{sourcepointer-8}{}\fi
\noindent\textit{Source reference.} Original reference 0CWF: sources/stacks-spaces-flat/pinned-original-source.tex, source lines 2030--2030. Prerequisite pointer invoked at proof line(s) 2279; it names a standard fact of the same chapter that the counted argument uses, and every step of the counted proof built on it is retained verbatim in the delivered body.\par\medskip
\ifdefined\hypertarget\hypertarget{sourcepointer-9}{}\fi
\noindent\textit{Source reference.} Original reference 00PP: sources/stacks-spaces-flat/pinned-original-source.tex, source lines 525--525. Prerequisite pointer invoked at proof line(s) 2221; it names a standard fact of the same chapter that the counted argument uses, and every step of the counted proof built on it is retained verbatim in the delivered body.\par\medskip
\ifdefined\hypertarget\hypertarget{sourcepointer-10}{}\fi
\noindent\textit{Source reference.} Original lemma 02W8: sources/stacks-chow/tags.txt, source lines 3795--3795. Prerequisite pointer invoked at proof line(s) 2237; it names a standard fact of a sibling chapter that the counted argument uses, and every step of the counted proof built on it is retained verbatim in the delivered body.\par\medskip
\ifdefined\hypertarget\hypertarget{sourcepointer-11}{}\fi
\noindent\textit{Source reference.} Original lemma 05LW: sources/stacks-chow/tags.txt, source lines 7250--7250. Prerequisite pointer invoked at proof line(s) 2292; it names a standard fact of a sibling chapter that the counted argument uses, and every step of the counted proof built on it is retained verbatim in the delivered body.\par\medskip
\ifdefined\hypertarget\hypertarget{sourcepointer-12}{}\fi
\noindent\textit{Source reference.} Original lemma 0CT8: sources/stacks-chow/tags.txt, source lines 16058--16058. Prerequisite pointer invoked at proof line(s) 2233; it names a standard fact of a sibling chapter that the counted argument uses, and every step of the counted proof built on it is retained verbatim in the delivered body.\par\medskip
\ifdefined\hypertarget\hypertarget{sourcepointer-13}{}\fi
\noindent\textit{Source reference.} Original reference 05MT: sources/stacks-spaces-flat/pinned-original-source.tex, source lines 2085--2085. Prerequisite pointer invoked at proof line(s) 2197; it names a standard fact of the same chapter that the counted argument uses, and every step of the counted proof built on it is retained verbatim in the delivered body.\par\medskip
\ifdefined\hypertarget\hypertarget{sourcepointer-14}{}\fi
\noindent\textit{Source reference.} Original lemma 0CVV: sources/stacks-chow/tags.txt, source lines 16150--16150. Prerequisite pointer invoked at proof line(s) 2289; it names a standard fact of a sibling chapter that the counted argument uses, and every step of the counted proof built on it is retained verbatim in the delivered body.\par\medskip
\begin{theorem}
\label{theorem-flat-dimension-n-representable}
In Situation \ref{situation-flat-dimension-n}.
Assume moreover that $f$ is of finite presentation, that
$\mathcal{F}$ is an $\mathcal{O}_X$-module of finite presentation,
and that $\mathcal{F}$ is pure relative to $Y$.
Then $F_n$ is an algebraic space and
$F_n \to Y$ is a monomorphism of finite presentation.
\end{theorem}

\begin{proof}
The functor $F_n$ is a sheaf for the fppf topology by
Lemma \ref{lemma-flat-dimension-n}.
Since $F_n \to Y$ is a monomorphism of sheaves on
$(\Sch/S)_{fppf}$ we see that $\Delta : F_n \to F_n \times F_n$
is the pullback of the diagonal $\Delta_Y : Y \to Y \times_S Y$.
Hence the representability of $\Delta_Y$ implies the same
thing for $F_n$. Therefore it suffices to prove that
there exists a scheme $W$ over $S$ and a surjective \'etale morphism
$W \to F_n$.

\medskip\noindent
To construct $W \to F_n$ choose an \'etale covering $\{Y_i \to Y\}$
with $Y_i$ a scheme. Let $X_i = X \times_Y Y_i$ and let
$\mathcal{F}_i$ be the pullback of $\mathcal{F}$ to $X_i$.
Then $\mathcal{F}_i$ is pure relative to $Y_i$ either by definition
or by Lemma \ref{lemma-quasi-finite-base-change}.
The other assumptions of the theorem are preserved as well.
Finally, the restriction of $F_n$ to $Y_i$ is
the functor $F_n$ corresponding to $X_i \to Y_i$ and $\mathcal{F}_i$.
Hence it suffices to show: Given $\mathcal{F}$ and $f : X \to Y$
as in the statement of the theorem where $Y$ is a scheme, the
functor $F_n$ is representable by a scheme $Z_n$ and
$Z_n \to Y$ is a monomorphism of finite presentation.

\medskip\noindent
Observe that a monomorphism of finite presentation is
separated and quasi-finite (Morphisms, Lemma
\ref{morphisms-lemma-monomorphism-loc-finite-type-loc-quasi-finite}).
Hence combining
Descent, Lemma \ref{descent-lemma-descent-data-sheaves},
More on Morphisms, Lemma
\ref{more-morphisms-lemma-separated-locally-quasi-finite-morphisms-fppf-descend}
, and
Descent, Lemmas \ref{descent-lemma-descending-property-monomorphism} and
\ref{descent-lemma-descending-property-finite-presentation}
we see that the question is local for the \'etale topology on $Y$.

\medskip\noindent
In particular the situation is local for the Zariski topology on $Y$
and we may assume that $Y$ is affine. In this case the dimension of the
fibres of $f$ is bounded above, hence we see that $F_n$ is representable
for $n$ large enough. Thus we may use descending induction on $n$.
Suppose that we know $F_{n + 1}$ is representable by a monomorphism
$Z_{n + 1} \to Y$ of finite presentation. Consider the base change
$X_{n + 1} = Z_{n + 1} \times_Y X$ and the pullback $\mathcal{F}_{n + 1}$
of $\mathcal{F}$ to $X_{n + 1}$. The morphism $Z_{n + 1} \to Y$ is
quasi-finite as it is a monomorphism of finite presentation, hence
Lemma \ref{lemma-quasi-finite-base-change}
implies that $\mathcal{F}_{n + 1}$ is pure relative to $Z_{n + 1}$.
Since $F_n$ is a subfunctor of $F_{n + 1}$ we conclude that in order
to prove the result for $F_n$ it suffices to prove the result for the
corresponding functor for the situation
$\mathcal{F}_{n + 1}/X_{n + 1}/Z_{n + 1}$.
In this way we reduce to proving the result for $F_n$ in case
$Y_{n + 1} = Y$, i.e., we may assume that $\mathcal{F}$ is flat
in dimensions $\geq n + 1$ over $Y$.

\medskip\noindent
Fix $n$ and assume $\mathcal{F}$ is flat in dimensions $\geq n + 1$
over the affine scheme $Y$.
To finish the proof we have to show that $F_n$ is representable
by a monomorphism $Z_n \to S$ of finite presentation.
Since the question is local in the \'etale topology on $Y$ it suffices to
show that for every $y \in Y$ there exists an \'etale neighbourhood
$(Y', y') \to (Y, y)$ such that the result holds after base change to $Y'$.
Thus by
Lemma \ref{lemma-existence-complete}
we may assume there exist \'etale morphisms $h_j : W_j \to X$,
$j = 1, \ldots, m$ such that for each $j$ there exists a complete
d\'evissage of $\mathcal{F}_j/W_j/Y$ over $y$,
where $\mathcal{F}_j$ is the pullback of $\mathcal{F}$ to $W_j$
and such that $|X_y| \subset \bigcup h_j(W_j)$. Since
$h_j$ is \'etale, by
Lemma \ref{lemma-pre-flat-dimension-n}
the sheaves $\mathcal{F}_j$ are still flat over in
dimensions $\geq n + 1$ over $Y$.
Set $W = \bigcup h_j(W_j)$, which is a quasi-compact open of $X$.
As $\mathcal{F}$ is pure along $X_y$ we see that
$$
E = \{t \in |Y| : \text{Ass}_{X_t}(\mathcal{F}_t) \subset W \}.
$$
contains all generalizations of $y$. By
Divisors on Spaces,
Lemma \ref{spaces-divisors-lemma-relative-assassin-constructible}
$E$ is a constructible subset of $Y$. We have seen that
$\Spec(\mathcal{O}_{Y, y}) \subset E$. By
Morphisms, Lemma \ref{morphisms-lemma-constructible-containing-open}
we see that $E$ contains an open neighbourhood of $y$.
Hence after shrinking $Y$ we may assume that $E = Y$.
It follows from
Lemma \ref{lemma-localize-flat-dimension-n}
that it suffices to prove the lemma for the functor $F_n$ associated to
$X = \coprod W_j$ and $\mathcal{F} = \coprod \mathcal{F}_j$.
If $F_{j, n}$ denotes the functor for $W_j \to Y$ and the sheaf
$\mathcal{F}_j$ we see that $F_n = \prod F_{j, n}$. Hence it suffices
to prove each $F_{j, n}$ is representable by some monomorphism
$Z_{j, n} \to Y$ of finite presentation, since then
$$
Z_n = Z_{1, n} \times_Y \ldots \times_Y Z_{m, n}
$$
Thus we have reduced the theorem to the special case handled in
More on Flatness, Lemma \ref{flat-lemma-flat-dimension-n-representable}.
\end{proof}

\end{document}
